\section{A continuous two-preparation marked family}
\label{sec:two-family}
The preceding saddle persists on an explicit interval of target laws,
with exact rather than perturbative values. Use the target $Q_\gamma$
from Section~\ref{sec:marked-exact} and write $U_2(\gamma)=U_2(b,Q_\gamma)$.

\begin{theorem}[A two-preparation saddle family]\label{thm:two-family}
For $6/25\le\gamma\le13/50$, let $r_\gamma$ be the unique root in
$(1/3,7/20)$ of
\begin{equation}\label{eq:family-cubic}
 f_\gamma(r)=r^3+(13-2\gamma)r^2+(3-12\gamma)r-(1+2\gamma)=0.
\end{equation}
Set
\[
 D_\gamma(r)=3r^2+(26-4\gamma)r+3-12\gamma,\qquad
 t_\gamma=-\frac{2\{3r_\gamma^2+(8\gamma-14)r_\gamma+3\}}
                         {D_\gamma(r_\gamma)}.
\]
Then $0<t_\gamma<1$, and
\begin{equation}\label{eq:two-family-value}
 U_2(\gamma)=
 \frac{-r_\gamma^3+(7-4\gamma)r_\gamma^2-3r_\gamma+13+4\gamma}{32}.
\end{equation}
The response \eqref{eq:two-response}, with $\tau=t_\gamma$, is optimal.
A least-favorable prior gives equal masses to the two diagonal parameters
$p=q=(1\pm\sqrt{r_\gamma})/2$. The response has the same five-event,
peak-twelve atomic realization. At $\gamma=1/4$ these formulas reduce to
Theorem~\ref{thm:two-saddle}.
\end{theorem}
\begin{proof}
For brevity write $r=r_\gamma$, $D=D_\gamma(r)$ and $t=t_\gamma$.
The derivative $D_\gamma$ is positive on the rectangle
$\gamma\in[6/25,13/50]$, $r\in[1/3,7/20]$. For example its lower bound
is $1/3+(26-52/50)/3-3/25>8$. The polynomial $f_\gamma$ decreases
with $\gamma$ at a fixed positive argument. Direct rational evaluations
give $f_{6/25}(1/3)<0$ and $f_{13/50}(7/20)>0$. Monotonicity therefore
gives the root and its uniqueness in this interval. The numerator of
$t$ is $-16\gamma r-6r^2+28r-6$. It is bounded below by
$28/3-6-16(13/50)(7/20)-6(7/20)^2>0$, and above by
$28(7/20)-6<4<D$. Thus $0<t<1$.

Let $H,J$ be the polynomials from \eqref{eq:poly} and
\eqref{eq:J-two}. Changing the target masses linearly in $\gamma$ gives
\begin{align*}
 H_\gamma(u,v)&=H(u,v)+(\gamma-1/4)\frac{1-(u-v)^2}{8},\\
 J_\gamma(u,v)&=J(u,v)+(\gamma-1/4)
                    \frac{(u-v)^2+6u-2v+1}{32}.
\end{align*}
The response's full-parameter score is $G_\gamma=H_\gamma+tJ_\gamma$.
One has $J_\gamma(r,0)=-f_\gamma(r)/64=0$ and
$\partial_uG_\gamma(r,0)=0$. Algebraic division gives
\begin{equation}\label{eq:family-square}
 G_\gamma(u,0)-H_\gamma(r,0)=\frac{(u-r)^2 A_\gamma(r,u)}{16D},
\end{equation}
where
\begin{align*}
 A_\gamma(r,u)={}&24\gamma^2+(3\gamma-10)r^2+(6\gamma-20)ru\\
                 &+12\gamma r+6\gamma u-51\gamma+30.
\end{align*}
For $0\le u\le1$ this is at least
$30-10(7/20)^2-20(7/20)-51(13/50)>0$.
The transverse derivatives satisfy
\begin{align*}
 32\partial_vH_\gamma
    &=3(u-v)^2+(8\gamma-10)u+(6-8\gamma)v+15\ge5,\\
 64\partial_vJ_\gamma
    &=3(u-v)^2+(14-4\gamma)u+(4\gamma-2)v-4\gamma-1\ge-3
\end{align*}
on $0\le u,v\le1$. Consequently $\partial_vG_\gamma\ge7/64$.
Equation~\eqref{eq:family-square} proves a lower bound
$H_\gamma(r,0)$, with equality precisely at the proposed prior points.
This quantity is the right side of \eqref{eq:two-family-value}.

For the upper bound group the full training words by $S$, as in the
preceding theorem. Put $L(r)=r^2+6r+1$. Using $f_\gamma(r)=0$, the first
three rows of $256C$ are
\begin{align*}
 0:&\quad (0,\;-8\gamma L(r),\;32r(r+1)-8\gamma L(r),\;32r(r+1)),\\
 1:&\quad 8(r-1)\bigl(r^2+4r-1-2\gamma(r+1),\bigr.\\[-2pt]
   &\hspace{30mm}r^2+4r-1+2\gamma(r+1),\\[-2pt]
   &\hspace{30mm}\bigl.r^2-4r-1+2\gamma(r+1),\;
                       r^2-4r-1-2\gamma(r+1)\bigr),\\
 2:&\quad-12(r-1)^2(r-1-2\gamma,\;r-1+2\gamma,\;
                                     r-1+2\gamma,\;r-1-2\gamma).
\end{align*}
The other rows are their reversals. The sign pattern is again
\[
 (0,-,+,+),\quad(+,-,+,+),\quad(+,+,+,+),\quad
 (+,+,-,+),\quad(+,+,-,0).
\]
Here are explicit interval checks for the potentially small margins.
In row zero the third coordinate is bounded below by
$32(1/3)(4/3)-8(13/50)((7/20)^2+6(7/20)+1)>0$.
The first parenthesis of row one is at most
$(7/20)^2+4(7/20)-1-2(6/25)(4/3)<0$; the third is at most
$(7/20)^2-4/3-1+2(13/50)(27/20)<0$.
The second parenthesis is positive already from $r>1/3$ and $\gamma>0$,
and the fourth is negative. Finally $r-1+2\gamma\le-13/100<0$.
All off-diagonal validation coefficients are nonpositive.
Thus the stated response is a Bayes best response for this prior,
including its two zero coefficients. Its score at both supporting
parameters equals the lower bound. The full posterior-predictive dual
proves \eqref{eq:two-family-value}.
\end{proof}

This interval is a uniform certificate region for the displayed saddle,
not a claim that the same masks remain optimal for every mark bias.
The theorem supplies a continuous family at genuinely growing training
history $N=2$; it does not infer a recurrence for every $N$ from the
one-preparation formulas.
